% !TEX root = 00_main_aaai_2027.tex
\section{Scope and Limitations}
The construction deliberately includes the oracle and changes distribution shape without changing the mean.
It demonstrates nonidentification by scalar summaries, not the prevalence of this failure among natural judges.
The family split holds out synthetic parameter bands; it is not evidence of transfer across real tasks, solvers, raters or model families.
Independent binomial ratings omit rater heterogeneity and dependence.
The study trains no language model and uses no new provider calls.

One observed rating can score a conditional forecast, but learning a flexible conditional distribution requires data and shared structure.
Neither a high correlation nor low finite-sample proper loss establishes that an entire human population has been learned.
When only a conditional mean or ranking is required, scalar objectives may be appropriate for that narrower target.
When the distribution is required, report proper probability losses and keep association, absolute error and bias as complementary diagnostics.

\section{Conclusion}
Correlation alone does not identify a judge's conditional rating distribution.
Proper probability losses identify that target in the population under their assumptions; sparse data and optimization still determine empirical selection quality.
Our controlled study supports this distinction while preserving inconclusive historical evidence.
